% ============================================================
% 203_R_Operator_De_ch.tex
% Der Rekursionsoperator R in FFGFT
% Formale Übersetzung der Herleitung aus Dok. 180
% Johann Pascher — Mai 2026
% ============================================================

\section*{Abstract}

Der Rekursionsoperator $\mathcal{R}$ der FFGFT formalisiert die
fraktale Skalierungsstruktur im Modenraum. Dieses Dokument übersetzt
die in Dok.~180 physikalisch hergeleitete Konstante
$\xi_0 = 4/30000$ in die formale Sprache des Operators und prüft,
ob die Fixpunktgleichung $\mathcal{R}(\Phi_*) = \Phi_*$ diese
Eindeutigkeit reproduziert. Das ist nicht der Fall: In keiner Lesart
liefert die Fixpunktbedingung $\xi_0 = 4/30000$; der Wert bleibt
durch die drei Bedingungen aus Dok.~180 festgelegt.

%------------------------------------------------------------
\section{Ausgangspunkt: Die physikalische Herleitung (Dok.~180)}
%------------------------------------------------------------

Dok.~180 leitet $\xi_0 = 4/30000$ aus drei unabhängigen Bedingungen her:

\begin{enumerate}
  \item \textbf{Dimensionsbeitrag $10^{-4}$:}
    Vier Raumzeit-Dimensionen liefern $\xi \sim (10^{-1})^4 = 10^{-4}$.
  \item \textbf{Geometriebeitrag $4/3$:}
    Tetraeder-Symmetrie: $4$ Ecken / $3$ Kanten pro Ecke $= 4/3$.
  \item \textbf{Selbstkonsistenz:}
    Das e-p-$\mu$-System liefert $\kappa = 7$ als einzige
    konsistente Rekursionstiefe ($0{,}04\,\%$ Genauigkeit).
\end{enumerate}

Diese drei Bedingungen legen $\xi_0$ ohne freien Parameter fest.
Die Aufgabe dieses Dokuments: Übersetzung in die $\mathcal{R}$-Sprache.

%------------------------------------------------------------
\section{Definition des Rekursionsoperators}
%------------------------------------------------------------

\subsection{Diagonale Form im Modenraum}

Der Rekursionsoperator $\mathcal{R}$ wirkt auf den Hilbertraum
$\mathcal{H} = \bigoplus_{n,l,j}\mathbb{C}\cdot\mathbf{e}_{n,l,j}$
diagonal:
\begin{equation}
  \mathcal{R}(\mathbf{e}_{n,l,j})
  = \lambda_{n,l,j} \cdot \mathbf{e}_{n,l,j}
  \label{eq:R-diagonal}
\end{equation}

Die Eigenwerte $\lambda_{n,l,j}$ sind durch die Massenformel bestimmt:
\begin{equation}
  \lambda_{n,l,j} = \xi_0^{p(n,l,j)-1} \cdot K_{\text{frak}}
  \label{eq:eigenwerte}
\end{equation}

mit $K_{\text{frak}} \approx 0{,}986$ (fraktale Korrektur).

\subsection{Wirkung auf den Feldoperator}

Auf dem Moden-Operator $\Phi = \sum_i m_i P_i$ (Dok.~202) wirkt
$\mathcal{R}$ nicht als einheitliche Skalierungsabbildung
$\mathcal{R}(\Phi) = \xi_0 \cdot \Phi$. Die Zeit-Masse-Dualitäts-Symmetrie
lautet
\begin{equation}
  t \to \xi_0 \cdot t, \qquad m \to \xi_0^{-1} \cdot m,
  \qquad \Phi \to \Phi .
  \label{eq:tmd-symmetrie}
\end{equation}
Mit $\Phi = \sum_i m_i P_i$ und $m \to m/\xi_0$ ergäbe sich
$\Phi/\xi_0$, nicht $\xi_0\Phi$; zugleich steht in
\eqref{eq:tmd-symmetrie} $\Phi \to \Phi$, und die Eigenwerte
\eqref{eq:eigenwerte} $\xi_0^{p-1}K_{\text{frak}}$ sind modenabhängig,
also keine einheitliche Skalierung mit $\xi_0$. Diese drei Angaben
sind nicht miteinander vereinbar; eine einheitliche Skalierung von
$\Phi$ unter $\mathcal{R}$ ist damit nicht festgelegt.

%------------------------------------------------------------
\section{Die Fixpunktgleichung}
%------------------------------------------------------------

\subsection{Definition des Fixpunkts}

Der Fixpunkt $\Phi_*$ ist der Operator, der unter $\mathcal{R}$
invariant ist:
\begin{equation}
  \mathcal{R}(\Phi_*) = \Phi_*
  \label{eq:fixpunkt}
\end{equation}

Eine einheitliche Skalierung $\mathcal{R}(\Phi) = s\,\Phi$ mit
$s = \xi_0$ oder $s = 1/\xi_0$ hätte nur diesen Fixpunkt:
\begin{equation}
  s \cdot \Phi_* = \Phi_*
  \quad \Rightarrow \quad
  (s - 1)\,\Phi_* = 0
\end{equation}

Da $s \neq 1$, folgt $\Phi_* = 0$ — ein trivialer Fixpunkt.

\subsection{Der nicht-triviale Fixpunkt durch Selbstkonsistenz}

Der physikalisch relevante Fixpunkt entsteht durch die
Selbstkonsistenzbedingung der Massenformel. Der Operator
$\Phi_* = \sum_i m_i^{(*)} P_i$ ist ein Fixpunkt von $\mathcal{R}$
wenn gilt:
\begin{equation}
  m_i^{(*)} = r_i \cdot \xi_0^{p_i} \cdot v
  \quad \text{und} \quad
  \mathcal{R}(m_i^{(*)}) = m_i^{(*)}
  \label{eq:selbstkonsistenz}
\end{equation}

Die zweite Bedingung bedeutet:
\begin{equation}
  \xi_0^{p_i - 1} \cdot K_{\text{frak}} \cdot m_i^{(*)}
  = m_i^{(*)}
  \quad \Rightarrow \quad
  \xi_0^{p_i - 1} \cdot K_{\text{frak}} = 1
  \label{eq:fixpunkt-bedingung}
\end{equation}

\subsection{Auflösung nach \texorpdfstring{$\xi_0$}{xi0}}

Aus \eqref{eq:fixpunkt-bedingung} folgt für jede Mode $i$:
\begin{equation}
  \xi_0 = K_{\text{frak}}^{-1/(p_i-1)}
  \label{eq:xi0-aus-fixpunkt}
\end{equation}

Damit $\xi_0$ \textbf{universell} (modunabhängig) ist, müsste diese
Gleichung für alle Moden gleichzeitig erfüllbar sein. Mit einem
universellen $K_{\text{frak}} \approx 0{,}986$ liefert sie jedoch
weder einen gemeinsamen Wert noch $\xi_0 = 4/30000$ (Schritt~3 unten).

%------------------------------------------------------------
\section{Verbindung zu den drei Bedingungen aus Dok.~180}
%------------------------------------------------------------

\subsection{Schritt 1: Dimensionsbeitrag}

Der Rekursionsoperator skaliert Längen mit $\xi_0$. In $D=4$
Raumzeit-Dimensionen skaliert das 4D-Volumen mit $\xi_0^4$:
\begin{equation}
  \mathcal{R}^{(4D)}(V_4) = \xi_0^4 \cdot V_4
\end{equation}

Für den Massenoperator, der mit $\ell^{-1}$ skaliert:
\begin{equation}
  \xi_0^{\text{dim}} = (10^{-1})^4 = 10^{-4}
\end{equation}

\subsection{Schritt 2: Geometriefaktor}

Die Tetraeder-Symmetrie des 4D-Torus liefert den geometrischen Faktor
aus der Topologie der Flächen:

\begin{center}
{\footnotesize
\adjustbox{max width=\textwidth}{%
\begin{tabular}{lcl}
\toprule
\textbf{Geometrie} & \textbf{Wert} & \textbf{Herkunft} \\
\midrule
Ecken des Tetraeders & 4 & Punktsymmetrie \\
Kanten pro Ecke & 3 & Raumeinbettung \\
Geometriefaktor $\xi_0^{\text{geo}}$ & $4/3$ & Verhältnis \\
\bottomrule
\end{tabular}%
}}
\end{center}

\subsection{Schritt 3: Selbstkonsistenz}

Die Fixpunktbedingung \eqref{eq:xi0-aus-fixpunkt} mit den Leptonen-Exponenten
$p_e = 3/2$, $p_\mu = 1$, $p_\tau = 2/3$ dient als Konsistenzprüfung.
Für das Elektron ($p_e - 1 = 1/2$):
\begin{equation}
  \xi_0 = K_{\text{frak}}^{-1/(3/2 - 1)}
  = K_{\text{frak}}^{-2}
  \approx (0{,}986)^{-2} \approx 1{,}03
\end{equation}
Das liegt weit von $\xi_0 = 4/30000 \approx 1{,}333 \cdot 10^{-4}$
entfernt; auch die Lesart $K_{\text{frak}}^2 = (0{,}986)^2 = 0{,}972$
trifft den Wert nicht. In keiner Lesart liefert die Fixpunktbedingung
$\xi_0 = 4/30000$.

Festgelegt wird $\xi_0$ daher allein über die drei Bedingungen aus
Dok.~180, mit dem e-p-$\mu$-System: der Exponent $\kappa = 7$ ist die
einzige ganzzahlige Rekursionstiefe, die alle drei Bedingungen
gleichzeitig erfüllt.

%------------------------------------------------------------
\section{Die Rekursionsformel}
%------------------------------------------------------------

Der Rekursionsoperator erzeugt die Massenhierarchie durch wiederholte
Anwendung:
\begin{equation}
  \mathcal{R}^{(k)}(\mathbf{e}_{n,l,j})
  = \xi_0^{k(p_{n,l,j}-1)} \cdot K_{\text{frak}}^k
  \cdot \mathbf{e}_{n,l,j}
  \label{eq:rekursion-k}
\end{equation}

Für $k = \kappa = 7$ (Rekursionstiefe aus Dok.~180):
\begin{equation}
  \mathcal{R}^{(7)}(\mathbf{e}_e) = \xi_0^{7(3/2-1)} \cdot K_{\text{frak}}^7
  \cdot \mathbf{e}_e
  = \xi_0^{7/2} \cdot K_{\text{frak}}^7 \cdot \mathbf{e}_e
\end{equation}

Die Selbstkonsistenzbedingung verlangt, dass dieser Ausdruck mit der
bekannten Elektronenmasse konsistent ist. Den Wert $\xi_0 = 4/30000$
fixiert er nicht; dieser wird aus Dok.~180 übernommen.

%------------------------------------------------------------
\section{Zusammenfassung: R als formaler Rahmen}
%------------------------------------------------------------

\begin{center}
{\footnotesize
\adjustbox{max width=\textwidth}{%
\begin{tabular}{p{4.5cm}p{6cm}}
\toprule
\textbf{Physikalische Herleitung} (Dok.~180) &
\textbf{Operator-Sprache} (Dok.~203) \\
\midrule
$10^{-4}$ aus 4D-Raumzeit &
  $\mathcal{R}^{(4D)}$ skaliert mit $\xi_0^4$ \\[4pt]
$4/3$ aus Tetraeder-Geometrie &
  Geometriefaktor in $K_{\text{frak}}$ \\[4pt]
$\kappa = 7$ aus e-p-$\mu$-System &
  $\mathcal{R}^{(7)}(\Phi_*) = \Phi_*$ Fixpunktbedingung \\[4pt]
$\xi_0 = 4/30000$ eindeutig &
  Übernommen; die Fixpunktbedingung liefert ihn nicht \\
\bottomrule
\end{tabular}%
}
}
\end{center}

Der Rekursionsoperator $\mathcal{R}$ ist damit kein neues Postulat,
sondern die formale Darstellung der in Dok.~180 abgeleiteten
Geometrie. Die Aufgabe aus Dok.~202, Problem~1 (Eigenwertform von
$\mathcal{R}$) ist damit formuliert; $\xi_0$ selbst wird nicht aus der
Fixpunktbedingung gewonnen, sondern aus Dok.~180 übernommen.

%------------------------------------------------------------
\section{Hinweis: Scheinbarer Widerspruch durch Kategorieverletzung}
%------------------------------------------------------------

Die Fixpunktbedingung $\mathcal{R}(\Phi_*) = \Phi_*$ legt nahe,
$\lambda_i = 1$ für jede Mode $i$ zu fordern. Einsetzen von
\eqref{eq:eigenwerte} liefert:
\[
  K_{\text{frak}}(e) = \xi_0^{-1/2} \approx 86{,}6, \quad
  K_{\text{frak}}(\mu) = 1, \quad
  K_{\text{frak}}(\tau) = \xi_0^{1/3} \approx 0{,}051
\]

Diese Werte sind alle verschieden — ein scheinbarer Widerspruch
zu einem universellen $K_{\text{frak}} \approx 0{,}986$.

\textbf{Aber:} Das ist \textbf{dieselbe Kategorieverletzung}
wie beim Nichtabschluss-Theorem (Dok.~193). Man vergleicht hier
Ausdrücke aus verschiedenen Modenräumen — genau so wie
$r_e \cdot r_\mu$ eine bedeutungslose Größe ist, weil $r_e$ und
$r_\mu$ Etiketten verschiedener algebraischer Räume sind,
ist der Vergleich $K_{\text{frak}}(e) = K_{\text{frak}}(\mu)$
unzulässig.

Die Fixpunktbedingung $\mathcal{R}(\Phi_*) = \Phi_*$ ist eine
\textbf{globale} Bedingung an den Feldoperator als Ganzes.
Sie ist \textbf{nicht} äquivalent zu der Forderung
$\lambda_i = 1$ für jedes $i$ separat, da die Projektoren
$P_i$ orthogonal sind und keine gemeinsame Skalierung erzwingen.

$K_{\text{frak}} \approx 0{,}986$ ist eine einzige geometrische
Konstante des fraktalen Torus (Dok.~180). Ihr Wert $1-100\,\xi_0$
ist über den Vergleich mit $\alpha$ bestätigt; die Form des fraktalen
Korrekturfaktors ist offen.


%------------------------------------------------------------
\section*{Zeichenerklärung}
%------------------------------------------------------------

\begin{center}
{\footnotesize
\resizebox{\textwidth}{!}{\begin{tabular}{lp{9cm}}
\toprule
\textbf{Symbol} & \textbf{Bedeutung} \\
\midrule
$\mathcal{R}$ & Rekursionsoperator im Modenraum \\
$\lambda_{n,l,j}$ & Eigenwert von $\mathcal{R}$ für Mode $(n,l,j)$ \\
$\lambda_{n,l,j} = \xi_0^{p-1}\cdot K_{\text{frak}}$ & Explizite Eigenwertformel \\
$K_{\text{frak}} \approx 0{,}986$ & Fraktale Korrekturkonstante (universal, Dok.~180) \\
$\kappa = 7$ & Rekursionstiefe (aus e-p-$\mu$-Selbstkonsistenz) \\
$\Phi_* = \sum_i m_i P_i$ & Fixpunktzustand des Operators \\
$\mathcal{R}(\Phi_*) = \Phi_*$ & Fixpunktbedingung \\
$p(n,l,j)$ & Exponentfunktion aus Torusgeometrie \\
$\Delta p_{\mu\to\tau} = 1/3$ & Schritt in geom. Folge $p_{n+1}=(2/3)p_n$; $p_\mu\cdot(1-2/3)=1/3$ \\
$\xi_0 = 4/30000$ & Geometrische Grundkonstante (Dok.~180) \\
$P_i = |\mathbf{e}_i\rangle\langle\mathbf{e}_i|$ & Projektor auf Mode $i$ (orthogonal) \\
$m_i = r_i\cdot\xi_0^{p_i}\cdot v$ & FFGFT-Massenformel \\
\bottomrule
\end{tabular}}}
\end{center}

\medskip\noindent
\textbf{Verweise:} Dok.~006 ($r_i$, $p_i$), Dok.~051 (Spinstruktur),
Dok.~180 ($\xi_0$-Herleitung), Dok.~189 (Stufenstruktur),
Dok.~193 (Nichtabschluss-Theorem), Dok.~202 (Übersicht).
